\providecommand{\todo}[1]{\textbf{[TODO: #1]}}

บทนี้อธิบายระเบียบวิธีดำเนินการวิจัย การจัดทำชุดข้อมูลทดสอบ สถาปัตยกรรมการประมวลผลสัญญาณเสียงและถอดความ ตรรกะการสกัดข้อมูลทางคลินิก การพัฒนาระบบเว็บแอปพลิเคชัน และระบบสมการตัวชี้วัดสำหรับการประเมินผล

\section{การวางแผนวิจัย}

โครงงานดำเนินงานตามขั้นตอน 4 ระยะ ได้แก่ (1) การจัดทำชุดข้อมูลทดสอบและเฉลยมาตรฐาน (2) การพัฒนากระบวนการประมวลผลเสียงและเอนจินถอดความ PathoWhisper (3) การพัฒนาตัวสกัดข้อมูลทางคลินิกและระบบเว็บ และ (4) การประเมินประสิทธิภาพของระบบ ผังการทำงานแสดงดังภาพที่ \ref{fig:workflow}

\begin{figure}[H]
\centering
\small
\begin{tabular}{c}
\fbox{\parbox{0.82\textwidth}{\centering 1. รับสัญญาณเสียงจากไมโครโฟนผ่านเบราว์เซอร์ หรืออัปโหลดไฟล์เสียง}}\\[2pt]
$\downarrow$\\[2pt]
\fbox{\parbox{0.82\textwidth}{\centering 2. กรองสัญญาณเสียงด้วยตัวกรอง afftdn $\rightarrow$ ตัดช่วงเงียบด้วย silenceremove $\rightarrow$ สุ่มแปลงเป็น 16 kHz Mono}}\\[2pt]
$\downarrow$\\[2pt]
\fbox{\parbox{0.82\textwidth}{\centering 3. ถอดความเป็นข้อความด้วย Faster-Whisper Small (INT8) ร่วมกับข้อความนำคลังศัพท์ CAP}}\\[2pt]
$\downarrow$\\[2pt]
\fbox{\parbox{0.82\textwidth}{\centering 4. ปรับข้อความและสกัดข้อมูลลงแบบฟอร์ม 15 ฟิลด์ด้วยตัวสกัดแบบกำหนดแน่นอน}}\\[2pt]
$\downarrow$\\[2pt]
\fbox{\parbox{0.82\textwidth}{\centering 5. แสดงผลให้พยาธิแพทย์ตรวจทานและยืนยัน $\rightarrow$ บันทึกฐานข้อมูล $\rightarrow$ ส่งออกเอกสาร PDF}}
\end{tabular}
\caption{ผังขั้นตอนการทำงานของระบบ PathoWhisper}
\label{fig:workflow}
\end{figure}

\section{การเก็บรวบรวมและจัดทำชุดข้อมูล}

\subsection{ชุดข้อมูลทดสอบ}

ชุดข้อมูลทดสอบประกอบด้วยกรณีศึกษาจำนวน 1,000 กรณี แบ่งออกเป็น 10 หมวดหมู่ หมวดละ 100 กรณี เพื่อประเมินขีดความสามารถของระบบในสภาวะแวดล้อมที่หลากหลาย ดังแสดงในตารางที่ \ref{tab:categories}

\begin{table}[H]
\centering
\caption{หมวดหมู่ของชุดข้อมูลทดสอบทางคลินิกจำลอง}
\label{tab:categories}
\small
\begin{tabular}{|c|p{5.5cm}|p{6.8cm}|}
\hline
\textbf{ที่} & \textbf{หมวดหมู่การทดสอบ} & \textbf{ลักษณะเงื่อนไขทางคลินิกที่ทดสอบ} \\
\hline
1 & Standard Protocol & บรรยายข้อมูลเรียงตามลำดับแบบฟอร์ม \\
\hline
2 & Out-of-Order Reporting & บรรยายข้อมูลสลับลำดับ เช่น บอกขนาดก้อนก่อนสิ่งส่งตรวจ \\
\hline
3 & Self-Correction \& Hesitation & มีการลังเลและแก้ไขตัวเลขหรือมิติกลางประโยค \\
\hline
4 & Multi-Margin Complex Assessment & มีการระบุขอบตัดและระยะห่างหลายทิศทาง \\
\hline
5 & Heavy Axillary Lymph Nodes Staging & บรรยายจำนวนและขนาดต่อมน้ำเหลืองที่พบด้วยตาเปล่า \\
\hline
6 & Fibrocystic / Benign Negative Control & สิ่งส่งตรวจที่ไม่ใช่มะเร็ง (กลุ่มควบคุมเชิงลบ) \\
\hline
7 & High-Speed Rapid Speech & อัตราการพูดเร็วเกิน 200 คำต่อนาที \\
\hline
8 & Fume Hood Noise ($-10$ dB) & เสียงพูดที่ซ้อนทับเสียงพัดลมตู้ดูดควันที่ SNR $-10$ dB \\
\hline
9 & Ductal Carcinoma / Microinvasion & ศัพท์เฉพาะเกี่ยวกับรอยโรค DCIS และการลุกลามขนาดเล็กn \\
\hline
10 & Clinical Edge Cases \& Missing Fields & กรณีที่มีข้อมูลไม่ครบถ้วนหรือมีฟิลด์เว้นว่าง \\
\hline
\end{tabular}
\end{table}

\subsection{รายละเอียดการสร้างชุดข้อมูล}
คุณลักษณะของชุดข้อมูลทดสอบสรุปดังตารางที่ \ref{tab:dataset} ข้อมูลเสียงสังเคราะห์ขึ้นจากสคริปต์ข้อความที่จัดโครงสร้างตามเกณฑ์มาตรฐาน CAP โดยมี 11 ฟิลด์ที่มีค่าบวกทางคลินิก และ 4 ฟิลด์ควบคุมเชิงลบที่ไม่มีการระบุค่า

\begin{table}[H]
\centering
\caption{คุณลักษณะของชุดข้อมูลทดสอบ}
\label{tab:dataset}
\small
\begin{tabular}{|p{4.2cm}|p{8.8cm}|}
\hline
\textbf{คุณลักษณะ} & \textbf{รายละเอียด} \\
\hline
จำนวนกรณีทดสอบ & 1,000 กรณีศึกษา (10 หมวด $\times$ 100 กรณี) \\
\hline
ไฟล์เสียงและรูปแบบ & 1,000 ไฟล์ รูปแบบ MP3 ความยาวเฉลี่ย 21.48 วินาทีต่อไฟล์ \\
\hline
การสังเคราะห์เสียง & Microsoft Edge TTS ภาษาอังกฤษ 2 สัญญาณเสียง (en-US-ChristopherNeural และ en-US-JennyNeural) \\
\hline
การจำลองสัญญาณรบกวน & ซ้อนทับเสียงพัดลมตู้ดูดควันพยาธิวิทยาจริงที่ SNR $-10$ dB ด้วย FFmpeg \\
\hline
เฉลยมาตรฐาน (Ground Truth) & สกัดจากสคริปต์แม่แบบตาม CAP Protocol ครบ 15 ฟิลด์ \\
\hline
\end{tabular}
\end{table}

\section{การประมวลผลและการสกัดข้อมูล}

\subsection{การประมวลผลสัญญาณเสียงเบื้องต้น}
สัญญาณเสียงอินพุตจะถูกประมวลผลผ่านตัวกรองสเปกตรัม \texttt{afftdn} ของ FFmpeg โดยกำหนดค่าการลดทอนสัญญาณรบกวน (\texttt{nr}) เท่ากับ 12 dB และระดับเสียงรบกวนพื้นฐาน (\texttt{nf}) เท่ากับ $-50$ dB ตามด้วยตัวกรอง \texttt{silenceremove} เพื่อตัดช่วงสัญญาณเงียบที่มีระดับพลังงานต่ำกว่า $-40$ dB โดยตัดช่วงเริ่มต้นที่ยาวเกิน 0.1 วินาที และช่วงหยุดพักระหว่างประโยคที่ยาวเกิน 0.6 วินาที ก่อนแปลงสัญญาณเป็น 16 kHz Mono สำหรับเอนจินถอดความ

\subsection{เอนจินการถอดความ PathoWhisper}
ระบบใช้ Faster-Whisper รุ่น Small (CTranslate2) ประมวลผลแบบ INT8 บน CPU จำนวน 8 เธรด กำหนด Beam Size เท่ากับ 1 เพื่อลดเวลาหน่วง พร้อมทั้งกำหนดข้อความนำบริบท (\texttt{PATHOLOGY\_PROMPT}) ซึ่งรวบรวมคำศัพท์ทางการแพทย์ตามเกณฑ์ CAP เพื่อชี้นำการสะกดคำที่ถูกต้อง โดยเปรียบเทียบกับแบบจำลองพื้นฐาน Whisper Small ต้นฉบับ (FP32 บน PyTorch, Beam Size 5, ไม่ใช้ Prompt)

\subsection{การปรับข้อความและการสกัดข้อมูลทางคลินิก}
ข้อความถอดความจะเข้าสู่ตัวปรับข้อความ (\texttt{normalizer.py}) เพื่อจัดการคำพูดแก้ไขกลางประโยค โดยตรวจจับคำสัญญาณ เช่น \textit{sorry, wait, actually, no wait} ที่ตามหลังตัวเลขมิติ แล้วเลือกเก็บเฉพาะค่าชุดล่าสุด ตลอดจนจัดการกรณีคำพ้องเสียง เช่น คำว่า \textit{weight} ที่ตามด้วยมิติชุดใหม่ให้ถือเป็นการแก้ไขค่า ตัวสกัดข้อมูลประยุกต์ใช้นิพจน์ปรกติในการจัดสรรข้อมูลลงแบบฟอร์ม 15 ฟิลด์ โดยจับคู่ชุดมิติ 3 มิติเข้ากับขนาดสิ่งส่งตรวจหรือขนาดก้อนตามคำระบุบริบทที่อยู่ใกล้ที่สุด ทั้งนี้ระบบทำงานแบบกำหนดแน่นอนและไม่สร้างข้อมูลที่ปรากฏไม่ชัดเจนในข้อความ

\section{การออกแบบและพัฒนาระบบ}

ระบบพัฒนาขึ้นเป็นเว็บแอปพลิเคชันด้วยภาษา Python และเฟรมเวิร์ก Flask ให้บริการแบบโลคอลผ่านโพรโทคอล HTTPS ฝั่งหน้าบ้านรองรับการบันทึกเสียงสดและการแสดงผลแบบฟอร์มเพื่อให้ผู้ตรวจทานสามารถแก้ไขข้อมูลได้แบบเรียลไทม์ ฝั่งหลังบ้านบันทึกข้อมูลที่ผ่านการยืนยันลงฐานข้อมูล PostgreSQL 16 และมี SQLite ทำหน้าที่เป็นฐานข้อมูลสำรองฉุกเฉินแบบออฟไลน์ พร้อมทั้งมีโมดูลส่งออกรายงานสรุปในรูปแบบเอกสาร PDF และ DOCX ผ่าน ReportLab และ PyMuPDF ซอร์สโค้ดของระบบและสคริปต์ที่ใช้ประเมินผลเผยแพร่ในคลังโค้ดสาธารณะ \cite{pathowhisperrepo}

\section{การประเมินผลและระบบสมการตัวชี้วัด}

\subsection{ระบบสมการตัวชี้วัดประสิทธิภาพ}
\label{sec:metrics}

เพื่อให้การประเมินผลมีความรัดกุมตามหลักวิชาการ โครงงานกำหนดตัวชี้วัดประสิทธิภาพเชิงคณิตศาสตร์ ดังนี้:

\subsubsection{อัตราความผิดพลาดระดับคำ (Word Error Rate: WER)} ใช้วัดความแม่นยำในการแปลงเสียงเป็นข้อความเทียบกับข้อความอ้างอิงเฉลย: 
\begin{equation} 
\mathrm{WER} = \frac{S + D + I}{N} \times 100\% 
\label{eq:wer} 
\end{equation} 
โดยที่ $S$ คือจำนวนคำที่ถูกแทนที่ผิด (Substitutions), $D$ คือจำนวนคำที่ตกหล่น (Deletions), $I$ คือจำนวนคำที่แทรกเกิน (Insertions), และ $N$ คือจำนวนคำทั้งหมดในข้อความเฉลย โดยคำนวณภายใต้ 2 เกณฑ์ ได้แก่ เกณฑ์ A (การจัดมาตรฐานข้อความพื้นฐาน) และเกณฑ์ B (การจัดมาตรฐานรูปแบบมิติและเลขที่สิ่งส่งตรวจ) เกณฑ์ B กำหนดหลังพิจารณาผลของเกณฑ์ A จึงรายงานทั้งสองเกณฑ์

\subsubsection{ความแม่นยำ การระลึก และคะแนนเอฟวันระดับฟิลด์ (Precision, Recall, $F_1$)}
ประเมินความถูกต้องของการสกัดข้อมูลลงแบบฟอร์มสำหรับฟิลด์ที่ $k$:
\begin{equation}
\mathrm{Precision}_k = \frac{\mathrm{TP}_k}{\mathrm{TP}_k + \mathrm{FP}_k} 
\times 100\%
\label{eq:precision}
\end{equation}
\begin{equation}
\mathrm{Recall}_k = \frac{\mathrm{TP}_k}{\mathrm{TP}_k + \mathrm{FN}_k} \times 100\%
\label{eq:recall}
\end{equation}
\begin{equation}
F_{1, k} = \frac{2 \cdot \mathrm{Precision}_k \cdot \mathrm{Recall}_k}{\mathrm{Precision}_k + \mathrm{Recall}_k} = \frac{2\mathrm{TP}_k}{2\mathrm{TP}_k + \mathrm{FP}_k + \mathrm{FN}_k} \times 100\%
\label{eq:f1} 
\end{equation}
โดยที่ $\mathrm{TP}_k, \mathrm{FP}_k, \mathrm{FN}_k$ แทนจำนวน True Positive, False Positive และ False Negative ของฟิลด์ที่ $k$ ตามลำดับ กรณีที่ระบบสกัดค่าผิดไปจากเฉลยจะถูกนับเป็นทั้ง $\mathrm{FP}$ และ $\mathrm{FN}$ พร้อมกัน

\subsubsection{ค่าเฉลี่ยมาโครของคะแนนเอฟวัน (Macro-$F_1$)}
ประเมินขีดความสามารถเฉลี่ยของการสกัดข้อมูลครอบคลุมฟิลด์ที่มีค่าบวกในชุดทดสอบจำนวน $K = 11$ ฟิลด์:
\begin{equation}
\mathrm{Macro\text{-}}F_1 = \frac{1}{K} \sum_{k=1}^{K} F_{1, k}
\label{eq:macrof1}
\end{equation}

\subsubsection{ความสอดคล้องระดับช่องข้อมูล (Slot-Level Agreement) และความจำเพาะ (Specificity)}
วัดสัดส่วนความถูกต้องรวมทุกช่องข้อมูล ($M = 15$ ฟิลด์ จาก $N_{\text{cases}} = 1{,}000$ กรณี รวมทั้งสิ้น 15,000 ช่องข้อมูล):
\begin{equation}
\mathrm{Slot\text{-}Agreement} = \frac{\sum_{k=1}^{M} (\mathrm{TP}_k + \mathrm{TN}_k)}{M \cdot N_{\text{cases}}} \times 100\%
\label{eq:slot_agreement}
\end{equation}
และประเมินความจำเพาะสำหรับกลุ่มฟิลด์ควบคุมเชิงลบ 4 ฟิลด์ ($k \in \{\text{s6, s7, s8, s9}\}$):
\begin{equation}
\mathrm{Specificity}_k = \frac{\mathrm{TN}_k}{\mathrm{TN}_k + \mathrm{FP}_k} \times 100\%
\label{eq:specificity}
\end{equation}
โดย $\mathrm{TN}_k$ แทนจำนวนช่องข้อมูลที่เฉลยเว้นว่างและระบบเว้นว่างได้ถูกต้อง

\subsubsection{อัตราความผิดพลาดวิกฤตระดับเคส ($\mathrm{CER}_{\mathrm{case}}$) และความถูกต้องสมบูรณ์}
ประเมินสัดส่วนกรณีศึกษาที่มีข้อผิดพลาดเกิดขึ้นในฟิลด์วิกฤตทางคลินิกอย่างน้อย 1 ฟิลด์ จากกลุ่มฟิลด์วิกฤต 6 ฟิลด์หลัก ($C_{\text{crit}} = \{\text{s1, s3, s10\_infiltrative, s10\_inf\_dims, s11, s14}\}$):
\begin{equation}
\mathrm{CER}_{\mathrm{case}} = \frac{1}{N_{\text{cases}}} \sum_{i=1}^{N_{\text{cases}}} \mathbb{I}\left( \sum_{c \in C_{\text{crit}}} (\mathrm{FP}_{i, c} + \mathrm{FN}_{i, c}) > 0 \right) \times 100\%
\label{eq:cercase}
\end{equation}
\begin{equation}
\mathrm{Exactness}_{\mathrm{case}} = 100\% - \mathrm{CER}_{\mathrm{case}}
\label{eq:exactness}
\end{equation}
โดยที่ $\mathbb{I}(\cdot)$ คือฟังก์ชันบ่งชี้ (Indicator Function) ซึ่งมีค่าเท่ากับ 1 เมื่อพบข้อผิดพลาดในเคส และเท่ากับ 0 เมื่อเคสดังกล่าวถูกต้องสมบูรณ์

\subsubsection{อัตราความผิดพลาดของมโนทัศน์ทางคลินิก (Concept Error Rate: ConER)}
วัดอัตราความผิดพลาดเฉพาะคำสำคัญทางคลินิก (คำศัพท์แพทย์ ตัวเลข และหน่วยมิติ):
\begin{equation}
\mathrm{ConER} = \frac{S_{\text{clinical}} + D_{\text{clinical}}}{N_{\text{clinical}}} \times 100\%
\label{eq:coner}
\end{equation}
โดย $S_{\text{clinical}}$ และ $D_{\text{clinical}}$ คือจำนวนคำสำคัญทางคลินิกที่ถูกแทนที่ผิดและตกหล่น และ $N_{\text{clinical}}$ คือจำนวนคำสำคัญทางคลินิกทั้งหมดในเฉลย

\subsubsection{ชนิดของความผิดพลาดและความครอบคลุมของธงเตือน}
แยกข้อผิดพลาดของการสกัดเป็นการเว้นว่าง (Omission) ซึ่งแพทย์เห็นช่องว่างได้ และการกรอกค่าผิด (Commission) ซึ่งผู้ตรวจอาจไม่สังเกต โดยให้ $r_{i,k}$ คือค่าเฉลยและ $h_{i,k}$ คือค่าที่สกัดได้ของฟิลด์ $k$ ในกรณีที่ $i$ และ $\varnothing$ คือช่องว่าง:
\begin{equation}
\mathrm{Omission} = \sum_{i=1}^{N_{\text{cases}}} \sum_{k=1}^{M} \mathbb{I}\left( r_{i,k} \neq \varnothing \wedge h_{i,k} = \varnothing \right)
\label{eq:omission}
\end{equation}
\begin{equation}
\mathrm{Commission} = \sum_{i=1}^{N_{\text{cases}}} \sum_{k=1}^{M} \mathbb{I}\left( h_{i,k} \neq \varnothing \wedge h_{i,k} \neq r_{i,k} \right)
\label{eq:commission}
\end{equation}
ความครอบคลุมของธงเตือน (Flag Coverage) คือสัดส่วนของเคสที่ฟิลด์วิกฤตผิดและระบบแสดงธงเตือนอย่างน้อยหนึ่งธง:
\begin{equation}
\mathrm{FlagCoverage} = \frac{N_{\text{flag}}}{N_{\text{err}}} \times 100\%
\label{eq:flagcov}
\end{equation}
โดยที่ $N_{\text{err}}$ คือจำนวนเคสที่ฟิลด์วิกฤตผิด และ $N_{\text{flag}}$ คือจำนวนเคสในกลุ่มนั้นที่ระบบแสดงธงเตือน

\subsubsection{เวลาประมวลผลและอัตราเร่ง (Latency and Speedup)}
เวลาประมวลผลคือค่าเฉลี่ยของเวลาต่อกรณี (วินาที) รวมการกรองเสียงและการถอดความของ PathoWhisper:
\begin{equation}
\mathrm{Latency} = \frac{1}{N_{\text{cases}}} \sum_{i=1}^{N_{\text{cases}}} t_i
\label{eq:latency}
\end{equation}
และอัตราเร่งคืออัตราส่วนระหว่างเวลาของแบบจำลองพื้นฐานกับของ PathoWhisper:
\begin{equation}
\mathrm{Speedup} = \frac{\mathrm{Latency}_{\mathrm{Baseline}}}{\mathrm{Latency}_{\mathrm{PathoWhisper}}}
\label{eq:speedup}
\end{equation}

\subsection{การวิเคราะห์ทางสถิติ}
การเปรียบเทียบผลระหว่างแบบจำลองดำเนินการแบบจับคู่รายกรณี (Paired Evaluation) ให้ $m_i^{\text{base}}$ และ $m_i^{\text{ours}}$ คือค่าตัวชี้วัด (เช่น WER หรือเวลา) ของกรณีที่ $i$ ในแบบจำลองพื้นฐานและใน PathoWhisper ผลต่างเฉลี่ยและการลดลงเชิงสัมพัทธ์คำนวณจาก:
\begin{equation}
\Delta = \frac{1}{N_{\text{cases}}} \sum_{i=1}^{N_{\text{cases}}} \left( m_i^{\text{base}} - m_i^{\text{ours}} \right)
\label{eq:delta}
\end{equation}
\begin{equation}
\mathrm{RelReduction} = \frac{\bar{m}^{\text{base}} - \bar{m}^{\text{ours}}}{\bar{m}^{\text{base}}} \times 100\%
\label{eq:relred}
\end{equation}
ช่วงความเชื่อมั่นร้อยละ 95 ของ $\Delta$ ได้จากวิธี Bootstrap จำนวน 5,000 รอบ และทดสอบนัยสำคัญด้วยการทดสอบวิลคอกซอน (Wilcoxon Signed-Rank Test) สำหรับค่าความผิดพลาดระดับคำและเวลา ส่วนความผิดพลาดระดับเคสใช้การทดสอบแมกนีมาร์ (McNemar's Exact Test) โดยให้ $b$ คือจำนวนเคสที่ baseline ผิดแต่ PathoWhisper ถูก และ $c$ คือจำนวนเคสที่ PathoWhisper ผิดแต่ baseline ถูก:
\begin{equation}
p = \min\left(1,\; 2 \sum_{j=0}^{\min(b,c)} \binom{b+c}{j} \left(\tfrac{1}{2}\right)^{b+c} \right)
\label{eq:mcnemar}
\end{equation}
